% ---------------------------------------------------------------
\section{Limitations and Future Work}
\label{sec:limitations}
% ---------------------------------------------------------------

The Safety, Liveness, and Accountability properties established in Section~\ref{sec:properties} are conditional guarantees. They hold under a defined set of structural assumptions about the host platform, the provisioning process, and the deployment context. This section makes those assumptions explicit, characterizes the operational consequences of their violation, and maps each gap to a directed research agenda. The purpose is to delimit the boundary of what the current specification guarantees — honestly — and to identify what remains open.

% ---------------------------------------------------------------
\subsection{Limitations}
% ---------------------------------------------------------------

\subsubsection{Human Response Latency.}
The Bailout Protocol's Liveness property guarantees eventual human contact within a bounded wall-clock interval $T_{\max}$, but $T_{\max}$ is a function of network topology and retry parameters — not of human cognitive response time. For operational domains with time constants measured in milliseconds — high-frequency trading, real-time process control, autonomous vehicle navigation — the time required for a human principal to receive, interpret, and issue a binding directive may exceed the system's operational time horizon by one to two orders of magnitude.

The protocol's prescribed failure mode under latency stress is well-defined: a node stalls in \texttt{ESCALATING} until the anchor responds or all pathways time out. This is the intended safety-preserving behavior. However, the operational cost of a stall is domain-dependent in ways the protocol does not currently reason about. In a precision agriculture context, a stall that delays a corrective action by minutes is acceptable; in a financial risk-interlock context, a stall that delays a margin-call response by even seconds may expose the system to catastrophic loss. The protocol treats the human anchor as a timing black box: there is no mechanism to distinguish whether delay reflects anchor overload, a structural mismatch between $T_{\max}$ and the anchor's operational availability, or a pathway failure.

The deeper limitation is that human response latency is provisioned statically at node setup and is not sensitive to the anchor's current operational context. $T_{\max}$ is a deployment parameter, not an adaptive runtime parameter. For latency-sensitive deployments, the provisioning burden falls on the system designer to establish empirically grounded latency budgets and encode them as enforceable policy. In deployment contexts where human response time is measured in seconds rather than minutes, static $T_{\max}$ provisioning may produce operational deadlocks that the protocol cannot distinguish from adversarial suppression.

\textit{Directed research:} Empirical characterization of human response latency distributions across deployment contexts to establish statistically grounded, domain-specific latency budgets. Integration of these budgets as enforceable policy parameters at node provisioning, with runtime monitoring of anchor responsiveness. Investigation of graded escalation responses: a protocol that requests a rapid \texttt{ACK} within seconds followed by a detailed directive within minutes would better match the operational reality of human decision-making under variable load.

\subsubsection{Adversarial Conditions.}
The Safety property requires that human-contact pathways cannot be permanently closed without recorded human authorization. This guarantee holds under the assumption that the host platform provides independent scheduling priority to the Bailout Monitor (Assumption A1.1). In an adversarial environment — particularly one in which the host platform itself may be compromised — this assumption may not hold. A sophisticated attacker controlling the execution substrate can modify thread priorities, preempt the monitor thread, introduce race conditions into pathway-closure operations, or rewrite the Bailout Monitor's code before it executes.

The protocol's guarantees are platform guarantees. They rest on invariants that the host platform must supply: non-preemptible monitor scheduling, append-only persistence for the authorization ledger, and interruptible blocking calls for pathway management. The Bailout Protocol is designed for production AI deployments in which individual agents may behave unexpectedly but the host platform is trusted — the standard assumption for regulated industries. It is not designed for deployments in which the platform itself is an active adversary: national infrastructure control systems, military command-and-control environments, financial clearinghouse infrastructure.

The distinction matters for deployment scoping. The protocol correctly prevents a misbehaving agent from suppressing its own escalation; it does not correctly prevent a compromised host from suppressing the monitor that enforces the protocol. This gap is structural, not fixable through protocol refinement alone. Additionally, the \fact{} layer pre-commit hook is the single integrity assumption for the entire guarantee structure. If the pre-commit hook is degraded — silently permitting writes it should reject — the safety and accountability properties dissolve without an internal detection mechanism, since any such detector would itself become the integrity assumption.

\textit{Directed research:} Formal specification of the composition interface between the Bailout Protocol and substrate-integrity mechanisms. Investigation of whether the protocol's guarantees can be preserved under weaker platform assumptions, including partially compromised substrates. Exploration of hardware-attested execution environments (e.g., Intel SGX, ARM TrustZone) as a deployment target for the Bailout Monitor, eliminating the OS as a trust dependency.

\subsubsection{Scalability.}
The protocol is specified with a 1:1 relationship between each node and its human anchor. This model is appropriate for single-node deployments, but it creates a structural bottleneck in enterprise contexts where one anchor may supervise hundreds of concurrent nodes. Without pooling and load-balancing mechanisms, the human anchor becomes the rate-limiting step in a system whose design intent is to route exceptions away from centralized machine decision points.

Consider a deployment with $N$ active nodes, each operating in an environment that produces exceptions at rate $\lambda$. The aggregate exception rate at a single anchor $h(n)$ is $N \cdot \lambda$. For sufficiently large $N$, exceptions queue faster than they can be acknowledged, and the system's effective liveness degrades because escalations remain pending beyond their operational relevance window. The protocol's $T_{\human}$ timeout converts unacknowledged exceptions to \texttt{RESOLVED} with \texttt{unresolved} tags — this is the correct safety-preserving behavior — but it does not solve the throughput problem. It documents the stall rather than preventing it.

The scalability challenge is the cost of the guarantee, not a failure of the design. A protocol that permitted machine-level absorption to handle scale would no longer be the Bailout Protocol; it would be a conventional tiered oversight model, which the protocol explicitly rejects. The architectural commitment to compulsory human contact has a provisioning cost that must be acknowledged and addressed.

Additional scalability considerations arise from spawn tree depth. As the depth of the accountability chain grows through recursive spawning, the time to propagate an exception to the anchor grows proportionally. At sufficiently large depths, the accumulated propagation delay may exceed $T_{\max}$ even under benign conditions, producing spurious \texttt{PATHWAY\_EXHAUSTED} resolutions that are indistinguishable from adversarial suppression in the Forensic Ledger.

\textit{Directed research:} Formal specification of human anchor pool architectures, including load-balancing policies, prioritization schemes, and routing algorithms that preserve the no-machine-actor property. Extension of Safety and Liveness proofs to pooled anchor configurations with formally verified pathway-redundancy guarantees. Empirical evaluation of anchor pool performance and spawn-tree-depth interactions under simulated high-escalation-load conditions.

\subsubsection{Exploitation Risk.}
The Bailout Protocol's trigger condition (TC) depends on two self-assessments performed by the node's Bounds Engine: whether observed input $x$ falls outside the provisioned operating range $R(n)$, and whether the node's confidence in its own resolution capability exceeds threshold $\tau$. Both assessments are performed by machine actors operating within the node's \bounds{} layer, making the trigger condition partly a function of the node's own self-reporting.

This creates an exploitation surface. A sufficiently sophisticated or adversarially modified node could systematically misclassify inputs that should trigger escalation as falling within $R(n)$, or could inflate its confidence score to suppress warranted escalation, converting the trigger condition from a structural safeguard into a machine-actor-controllable gate. The protocol's immutability constraints ($\tau$ is provisioned at initialization and cannot be adjusted at runtime by any machine actor) limit the scope of this exploitation, but they do not eliminate it: a node compromised at provisioning time can select an appropriately high $\tau$ value before the immutability constraint takes effect.

The Forensic Ledger provides post-hoc detectability for systematic misclassification: an auditor can reconstruct the relationship between observed input conditions and triggered escalations for each node over rolling observation windows and flag statistically anomalous patterns. However, this detection is forensic, not preventive. The protocol does not currently include an active anomaly detector that intercepts misclassification before it produces an unrecorded escalation.

The protocol also assumes human anchors are competent, attentive, and not subject to social engineering. Once an exception reaches $h(n)$, the protocol provides diagnostic context and a propagation chain, but it does not authenticate urgency or severity beyond what is recorded in the Forensic Ledger. An adversarial agent could craft an exception that appears routine — carrying a low-severity tag and minimal diagnostic context — to cause $h(n)$ to issue a directive without full scrutiny. The immutable attribution record partially mitigates this: any directive issued under insufficient scrutiny is permanently attributed to the human anchor. However, this is a deterrent, not a structural prevention. Fully automated detection of social engineering is outside the scope of any accountability architecture; it must be addressed through training, organizational process, and deployment-specific policies.

\textit{Directed research:} Triggering-behavior anomaly detectors that audit the statistical relationship between observed exception conditions and triggered escalations per node over rolling observation windows. Entropy-based monitors: a node that observes conditions outside its operating range but triggers escalation at a rate statistically inconsistent with its observed condition distribution is a candidate for audit. Investigation of provenance attestation for operating-range assessments: requiring that each Bounds Engine self-assessment be counter-signed by an independent audit process, making post-hoc misclassification detectable through forensic reconstruction. Runtime attestation that the node's $R(n)$ definition has not been altered since provisioning.

% ---------------------------------------------------------------
\subsection{Future Work}
% ---------------------------------------------------------------

The following three research directions extend the Bailout Protocol's guarantees to broader deployment contexts and establish a principled research program beyond the current specification.

\medskip\noindent\textbf{Formal verification.} Full mechanical verification of the Safety, Liveness, and Accountability invariants using a formal verification tool — Coq, Isabelle/HOL, or TLA$^+$ — applied to the three-layer model and the Bailout Protocol state machine. The current paper establishes the invariants through proof sketches grounded in the deterministic transition relation; formal verification is a prerequisite for deployment in regulated industries where proof of correctness is a certification requirement. Of particular interest is whether the functional separation between \purpose{}, \bounds{}, and \fact{} layers admits a compositional verification argument: if each layer satisfies its specification, does the composition necessarily satisfy the protocol's guarantees? A compositional verification architecture would distribute the verification burden across layer implementations rather than centralizing it on the full system. The \fact{} layer pre-commit hook is the primary target: if the pre-commit hook is incorrect, the entire upstream accountability guarantee is void. Formal verification is not a refinement of the protocol — it is a prerequisite for safety-critical deployment.

\medskip\noindent\textbf{Adaptive timeout mechanisms.} Replacement of static timeout parameters ($T_{\bounds}$, $T_{\prop}$, $T_{\cap}$, $T_{\human}$) with adaptive mechanisms that adjust to anchor workload, pathway latency, and deployment context while preserving the Safety invariant. Correct provisioning of static timeout values requires domain expertise and operational data that may not be available at deployment time; the primary motivation for adaptive timeout tuning is empirical measurement of human response distributions, which is only available \emph{after} the system is operational. The adaptive protocol would maintain the guarantee that the anchor is always contacted within $T_{\max}$ while reducing unnecessary stalls in low-consequence, high-latency scenarios. On each escalation cycle, the protocol records the actual response time $t_{\mathrm{actual}}$. Over a rolling observation window, it computes a distribution and sets $T_{\bounds}$ to a high percentile of that distribution plus a safety margin. Adaptation must be slow enough that a single anomalous response does not immediately distort the timeout, and Safety is preserved by never adjusting $T_{\bounds}$ below the minimum viable threshold established at provisioning. The more sophisticated extension — human anchor load balancing — routes exceptions to the anchor with the lowest current queue depth rather than to the statically provisioned $h(n)$, addressing the scalability bottleneck while preserving attribution: the Forensic Ledger records which anchor received each exception, and routing is constrained to anchors with appropriate domain authority over the originating node.

\medskip\noindent\textbf{Distributed human anchor pools.} Formal specification and empirical evaluation of anchor pool architectures that sustain scalability under cascading failure conditions without compromising the no-machine-actor-absorbs property. A distributed human anchor pool is a set of principals $\{h_1, h_2, \ldots, h_k\}$ each with authority over a subset of the node population, where exceptions are routed to the pool member with current availability and appropriate domain authority.

Extending the protocol to a pool requires addressing three hard problems. \textbf{Domain authority:} which pool members are authorized to receive exceptions from which nodes — this must be established at provisioning and enforced by the \fact{} layer. \textbf{Accountability:} the Forensic Ledger must record which pool member received each exception so that attribution is preserved even under load balancing; routing decisions must be committed to the ledger as protocol events. \textbf{Pool membership integrity:} a compromised node cannot add a non-human actor to the pool to create a false human anchor — this requires a provisioning protocol with cryptographic attestation that each pool member is a verified natural person or an institutional actor with unambiguous human authority.

The pool architecture must satisfy three constraints derived directly from the correctness properties. \textbf{Safety:} no pool configuration may produce a state in which all human-contact pathways for a given node are permanently disabled; each node must maintain at least two anchor references and the pool management algorithm must detect anchor unavailability before it becomes pathway-closing. \textbf{Liveness:} routing must deliver escalations to available anchors within $T_{\max}$ even under high concurrent load, requiring load-aware routing rather than static assignment. \textbf{Accountability:} every escalation must be attributed to a specific human anchor, immutably recorded. The pool architecture must also define behavior for when all anchors in a pool are simultaneously overloaded — a load-shedding policy that preserves Safety while providing meaningful feedback about the overload condition.

% ---------------------------------------------------------------
\section{Conclusion}
\label{sec:conclusion}
% ---------------------------------------------------------------

The Bailout Protocol advances a single, architecturally precise claim: that any multi-agent system built on the \bxthree{} Framework must propagate every unresolved exception state to a named human principal, bypassing all intermediate machine actors, as a compulsory structural requirement rather than a runtime choice. This mandatory bailout property is not a feature that can be retrofitted to an existing agentic architecture through better error handling or improved escalation heuristics. It is a structural property that follows, as a matter of architectural necessity, from the functional separation between the \purpose{}, \bounds{}, and \fact{} layers.

The argument is direct. When the \bounds{} layer encounters a condition outside its provisioned operating range $R(n)$ that it cannot resolve, no machine actor has the authority to adjudicate that condition. The \bounds{} layer detects exception conditions but carries no judgment authority; the \fact{} layer enforces transition relations and records events but carries no intent authority. The decision must return to the \purpose{} layer, which carries intent, judgment, and final accountability. The Bailout Protocol is the mechanism that enforces this return — compulsorily, deterministically, and with full forensic attribution at every step.

This contribution is realized within the \bxthree{} Framework as its fifth named architectural pillar, complementing Loop Isolation, Recursive Spawning, Spatial Firewall, and Sandbox Gate. Together, these five pillars constitute the upstream accountability guarantee: that \bxthree{} systems \emph{fail upward into human consciousness, never downward into autonomous chaos}. The Bailout Protocol is the runtime expression of this guarantee for all exception conditions not resolved by the preceding four pillars — including conditions not anticipated at design time, which is precisely the operational reality that makes the protocol necessary.

The \bxthree{} Framework's role in this contribution is to provide the conceptual architecture within which the Bailout Protocol operates as a necessary consequence rather than an independent design choice. The functional separation of Purpose, Bounds, and Fact layers creates the structural conditions that make the bailout requirement compulsory. The Forensic Ledger maintains the immutable attribution record that makes every escalation auditable. The immutable parent-pointer established at node provisioning ensures that the escalation path is fixed at the architectural level and cannot be modified by any machine actor. These elements are not independent mechanisms that happen to cooperate; they are mutually defining structural facts whose necessity is established by the protocol's own specification.

The theoretical claim is foundational: that the upstream accountability guarantee of the \bxthree{} Framework is not a design aspiration but a logical consequence of the framework's architecture, and that the Bailout Protocol is the mechanism by which that consequence is enforced at runtime. The practical claim is equally clear: that autonomous systems operating without such a mechanism are structurally exposed to failure modes — invisible exception absorption, contingent escalation, and irreversible drift — that cannot be eliminated by incremental improvement to existing error-handling or oversight practices. They require a new architectural primitive. The Bailout Protocol is that primitive.

The limitations enumerated in Section~\ref{sec:limitations} are not refutations of the protocol's contribution; they are the honest boundary conditions that define where the contribution applies. Human response latency, adversarial conditions, scalability, and exploitation risk are constraints on the deployment context, not flaws in the protocol's logical design. They constitute the research agenda for the work that follows: formal verification to establish the correctness proofs rigorously; adaptive timeout to reduce provisioning burden and improve operational fit; and distributed human anchor pools to enable deployment at industrial scale while preserving the Safety, Liveness, and Accountability guarantees. The protocol's structural necessity within the \bxthree{} Framework is established by its specification. Its practical necessity in the field is established by the operational reality of autonomous systems that, without it, have no architectural path to human accountability when they need it most.